\documentclass[UTF8,a4paper,11pt]{ctexart}
\usepackage[margin=2.2cm]{geometry}
\usepackage{booktabs,longtable,array,enumitem,hyperref,xcolor}
\hypersetup{colorlinks=true,linkcolor=blue,urlcolor=blue}
\setlist{nosep}
\title{AI工具使用详情}
\author{2026年全国大学生数学建模竞赛 C 题项目组}
\date{}
\begin{document}
\maketitle
\tableofcontents
\newpage
\section{AI工具使用概况}
本项目研究含光伏、负荷、储能和外购电的微网调度，依次处理确定性日前计划、预测误差下的风险修正、0:00/6:00/12:00/18:00 光伏预报更新，以及波动电价扩展。团队将 AI 定位为候选方案讨论、代码辅助、约束审查、结果解释和论文整理工具；数学结构、实验、数值结果和最终结论由团队完成并核验。
\begin{longtable}{p{0.2\textwidth}p{0.2\textwidth}p{0.27\textwidth}p{0.23\textwidth}}
\toprule AI工具 & 模型/版本 & 使用目的 & 主要环节\\\midrule
DeepSeek & DeepSeek V4.1 & 题意拆解、模型比较、假设审查、结果解释和文字检查 & 题目分析、问题一至四方案论证、结果复核\\
Kimi & Kimi K3 & 长文本材料梳理、时序约束审查、结果对照和论文结构检查 & 题目分析、问题三事件时序、问题四信息边界\\
GLM & GLM 5.2 & Python 编程、Debug、接口和绘图检查 & 核心模块、问题一至四求解程序\\
\bottomrule
\end{longtable}
\section{各建模环节中的具体使用}
\subsection{问题一：确定性日前购电计划}
最终采用 10 分钟时段确定性线性规划，使用 HiGHS 求解器。约束包括购电非负、光伏利用、储能充放电功率、SOC 递推、SOC 上下限和日初日末状态。结果为全天购电量 59482.70 kWh、现金购电费 35126.95 元，日初和日末 SOC 均为 6000 kWh。DeepSeek V4.1 与 Kimi K3 帮助梳理功率到电量的转换、母线平衡和 SOC 递推；GLM 5.2 协助变量索引、约束矩阵、结果导出和绘图。团队重新推导约束，并以无储能方案、能量守恒式和小样本手算复核。
\subsection{问题二：风险修正与滚动执行}
最终采用 F1 因果预测、历史净负荷残差分位安全轨迹和滚动线性规划。F1 的负荷 MAE、光伏 MAE、净负荷 MAE 和 RMSE 分别为 176.80、153.40、273.35、389.97 kW。主配置 F1-W30-alpha0.9 的 2025 年 2 月至 12 月现金总费为 14917501.49 元，应急购电量为 156695.85 kWh。团队采纳统一回测和闭环费用评价，但未采用 LSTM/Transformer 作为主模型；GLM 5.2 协助跨日 SOC、残差窗口、分位数和滚动 LP 的代码检查。
\subsection{问题三：事件驱动合同调整}
最终采用 0:00 初始合同、6:00/12:00/18:00 三次调整的事件驱动合同 LP，事件之间使用滚动 LP 和槽内裁剪。已执行时段和已发生现金流冻结，最终合同相对原合同结算一次。全量更新方案现金总费 14968496.84 元，应急购电量 87488.76 kWh；相对无日内更新方案节省 618746.73 元。AI 用于审查预测发布时间、合同生效、真实 SOC 和结算时序；团队据此固定四个授权时点。
\subsection{问题四：波动电价扩展}
问题四区分预测电价和真实结算电价，比较四种价格预测基线，主配置采用近 30 日中位数。固定合同模式现金总费为 17258751.73 元；可调合同模式为 15769829.13 元，少 1488922.60 元，应急购电量由 425014.50 kWh 降至 86353.65 kWh。AI 帮助检查因果信息边界并建议同时报告预测误差、尖峰误差和闭环现金费用。
\section{主要提示方式和使用过程}
团队使用模型讨论、代码审查和结果审查三类提示。提示通常包含题目条件、数据规模、已有方案、约束、函数输入输出、报错或结果摘要，要求 AI 比较候选方法、检查边界和信息泄漏，而不是直接完成模型。团队先形成基本模型，再让 AI 做比较和反例审查；代码建议独立运行，数值结论回到原始数据和统一脚本核验。
\section{典型交互示例}
\subsection{预测模型选择}
提示：“数据为逐日、逐10分钟负荷和光伏，预测只能使用目标日前信息。请比较 F0/F1/F2 与 LSTM 类模型，说明评价指标、信息泄漏和过拟合风险，不要直接替我们选模型。” AI 建议先统一因果回测并同时观察净负荷误差、应急购电和现金费用。团队采纳回测，未采用 LSTM/Transformer 主方案；通过问题二误差表、逐日历史范围审计和年度闭环费用核验。
\subsection{风险分位参数}
提示：“请解释 alpha=0.8、0.9、0.95 的含义，并说明为何不能把它直接解释为整日可靠度。” AI 指出分位不等于整日联合覆盖率。团队修改后采纳风险框架，未照搬默认参数；通过 W/alpha、效率、容量、功率和终端价值敏感性核验。
\subsection{滚动 LP Debug}
提示：“充放电效率均为0.9，x/y为交流侧kWh，请检查 SOC 递推符号、变量索引、边界、首槽执行和应急电主动充电限制。” AI 提醒核对 $S_{t+1}=S_t+0.9x_t-y_t/0.9$。团队修改边界和槽内裁剪，并用小样本手算、母线平衡和全年门禁核验。
\subsection{问题三事件时序}
提示：“只允许0/6/12/18获得预报并调整剩余合同，请检查未来信息、已执行时段覆盖和重复结算。” AI 建议事件流、冻结已执行时段、一次性结算。团队采纳并通过事件抽查、八组消融和逐槽/逐日费用核对。
\subsection{波动电价评价}
提示：“说明价格 MAE、RMSE、尖峰 MAE 与全年现金费用可能不一致的原因，并设计完全价格信息参照。” AI 建议联合报告误差、应急量、现金费，并将完全信息作为不可实施基准。团队采纳评价框架，未把完全信息组当实际策略。
\subsection{结果解释修正}
AI 初稿使用“光伏增加导致费用下降”等强因果措辞。团队要求改为“在该策略和信息集下伴随”等条件性表述，并逐项对照消融和敏感性结果，删除无证据结论。
\section{AI输出的采纳、修改和推翻}
\begin{longtable}{p{0.25\textwidth}p{0.15\textwidth}p{0.16\textwidth}p{0.32\textwidth}}
\toprule AI输出/建议 & 对应环节 & 团队处理 & 原因与核验\\\midrule
统一回测 F0/F1/F2 & 问题二 & 直接采纳 & 逐日因果回放、MAE/RMSE、费用比较\\
残差分位安全轨迹 & 问题二 & 修改后采纳 & 分位不等于整日可靠度；做 W/alpha 敏感性\\
LSTM/Transformer 主模型 & 问题二 & 验证后推翻 & 数据规模、解释性和闭环效果不支持\\
滚动 LP 代码框架 & 问题二至四 & 修改后采纳 & 修正效率、边界、跨日 SOC 和裁剪\\
连续获得新光伏预报 & 问题三 & 否决 & 超出授权时点，造成信息泄漏\\
最低价格 MAE 直接选策略 & 问题四 & 未采纳 & 预测误差与调度费用不等价\\
“导致”等强因果表述 & 结果解释 & 修改后采纳 & 回放不足以证明因果\\
\bottomrule
\end{longtable}
最明确的推翻链条是：AI 提出复杂预测模型方向，团队在共同信息集下回测，因数据规模、解释性和闭环表现不支持而不采用。最明确的修改链条是：AI 识别滚动 LP 的效率和边界风险，团队结合题目约束重写代码并通过门禁。
\section{人工核验方式}
\subsection{数学模型核验}重新推导功率到电量、SOC 递推、母线平衡、应急费用和合同调整结算式；检查非负性、SOC 1200--10800 kWh、充放电互斥和应急电不得主动充电。问题一用日初日末 SOC 和全天能量守恒式独立校验。
\subsection{程序核验}人工阅读核心模块和四个求解程序，核对变量索引、10分钟步长、跨日 SOC、预测窗口和事件冻结；对边界和窗口末端进行小样本测试。全年结果通过统一门禁，费用加总误差为 0 或浮点舍入量级，且无回退。
\subsection{数据核验}检查附件 1 的 144 个时点、附件 2/4 的 365 日连续日期、附件 3 的四个发布时间和每次 24 个预测值，核对缺失、负值、重复日期和时间标签。预测回测只读取目标日前数据。
\subsection{模型效果核验}问题二联合比较 MAE、RMSE、应急量和现金总费；问题三用八组信息消融区分新预报价值与重算机会；问题四报告价格误差、合同模式费用和完全信息基准；另做效率、容量、功率、初始 SOC、终端价值和时间标签敏感性。
\subsection{最终结果核验}逐槽 CSV、逐日汇总、Excel 工作簿和论文表格互相对照，人工抽查 2025-03-20、06-21、09-23、12-21，并核对 Q3/Q4-3 每个目标槽只相对原合同结算一次。
\section{AI参与边界}
AI 参与思路扩展、候选模型比较、约束和信息时序审查、Python 代码辅助、Debug、参数分析建议、图表组织、结果解释和语言润色。团队负责合同决策层—物理执行层—现金结算层的统一框架、关键假设、核心推导、参数最终确定、数值实验、模型优劣判断和论文结论。AI 不替代团队完成最终模型。
\section{总结}
AI 提高了方案探索、代码编写、程序审查、结果整理和文字表达效率，也帮助发现时间标签、信息因果性、SOC 递推、重复结算和过强因果表述等风险。复杂预测模型、连续预报更新和单纯依据价格 MAE 选策略等建议均未被直接接受，部分建议经过实验后被修改或推翻。最终模型及结论来自团队自己的数学分析、程序计算、统一信息集比较、年度闭环回放、敏感性分析和人工核验。
\end{document}
